% (User input window)

\paragraph{跨接口终端结论：真实输入核的记忆可观测性、strictification 普适性与端点/熵率审计常数族}
本段把真实输入核、strictification（完全正可见实现）与端点审计三条链条中的可复核常数与普适性质，
压缩为可独立引用的结论条目；证明仅指向相应已完成的定理/命题/推论（或其可审计输出），不再重复推导细节。

\begin{conclusion}[真实输入核的 Parry 记忆位严格解耦：\texorpdfstring{$(p_x,p_y)$}{(px,py)} 在最大熵下给出闭式边缘并严格独立]\label{con:xi-terminal-real-input-parry-memory-decoupling}
在本质 \(20\) 状态真实输入核 \(M_{\mathrm{ess}}\) 的 Parry 最大熵测度下，输入记忆
\(m=(p_x,p_y)\in\{00,01,10,11\}\) 的边缘分布满足闭式
\[
\mathbb{P}(00)=\frac{3+\sqrt5}{10},\qquad
\mathbb{P}(01)=\mathbb{P}(10)=\frac15,\qquad
\mathbb{P}(11)=\frac{3-\sqrt5}{10}.
\]
并且两比特严格独立且同边缘：
\[
\mathbb{P}(p_x=1)=\mathbb{P}(p_y=1)=\frac{5-\sqrt5}{10},
\qquad
\mathbb{P}(p_x,p_y)=\mathbb{P}(p_x)\mathbb{P}(p_y).
\]
该解耦为精确代数事实（而非近似），可由同步态条件分布表的支撑结构反推强制得到（例如特定同步态仅在单一记忆态下出现）。
\end{conclusion}
\begin{proof}[证明]
见命题 \ref{prop:real-input-40-memory-marginal} 与推论 \ref{cor:real-input-40-memory-independence}（其推导中使用了表 \ref{tab:real_input_40_parry_internal_distribution} 的支撑结构）。证毕。
\end{proof}

\begin{conclusion}[同步态—记忆信道的硬判定分层与可观测性常数：互信息与 MAP 准确率出现 \texorpdfstring{$\sqrt5$}{sqrt5} 闭式]\label{con:xi-terminal-sync-memory-hard-decision-map-accuracy}
把同步态 \(S\in Q\) 视为对记忆 \(M=(p_x,p_y)\) 的观测信道。
在 \(M_{\mathrm{ess}}\) 的 Parry 最大熵测度下，存在一组硬判定同步态使后验退化为确定值：
\[
s\in\{01\text{-}1,010,11\text{-}1\}\Rightarrow \mathbb{P}(M=00\mid s)=1,
\qquad
s\in\{0\text{-}12,1\text{-}12\}\Rightarrow \mathbb{P}(M=11\mid s)=1,
\]
且这类硬判定同步态的总概率为
\[
\mathbb{P}(01\text{-}1)+\mathbb{P}(010)+\mathbb{P}(11\text{-}1)+\mathbb{P}(0\text{-}12)+\mathbb{P}(1\text{-}12)=\frac{3}{10}.
\]
此外，互信息为可复核常数
\(
I(S;M)\approx 0.812300
\) bits。
以 MAP 规则在线判别记忆（取使 \(\mathbb{P}(\cdot\mid s)\) 最大者），最优可达准确率满足闭式
\[
\mathrm{Acc}^*(M\leftarrow S)=\sum_s \mathbb{P}(s)\max_m \mathbb{P}(m\mid s)=\frac{3\sqrt5}{10}.
\]
对若干自然二值函数亦有闭式最优准确率：
\[
\mathrm{Acc}^*(p_x\leftarrow S)=\mathrm{Acc}^*(p_y\leftarrow S)=\frac45,
\qquad
\mathrm{Acc}^*(C\leftarrow S)=\frac{17+\sqrt5}{20},
\quad
C:=\ind_{\{M=11\}}.
\]
\end{conclusion}
\begin{proof}[证明]
硬判定分层与其总概率见注 \ref{rem:real-input-40-bayes-decode}；互信息常数见注 \ref{rem:real-input-40-mutual-info}；
MAP 准确率及其派生二值函数的闭式见命题 \ref{prop:real-input-40-map-accuracies}。证毕。
\end{proof}

\begin{conclusion}[strictification 的普适表征：\texorpdfstring{$G^{\mathrm{BC}}_m$}{GBCm} 是“完全正可见实现存在性”的最大可见群商]\label{con:xi-terminal-strictification-maximal-cp-visible-quotient}
固定分辨率 \(m\) 并令 \(\mathcal{A}_m:=\CC[G^{\mathrm{BC}}_m]\)。
若可见扭曲通道族 \(\{F_{\overline\chi}\}_{\overline\chi\in\widehat{G^{\mathrm{BC}}_m}}\) 可组装为 \(\mathcal{A}_m\)-值解析函数 \(F_m:\mathbb{D}\to\mathcal{A}_m\)，且
\begin{enumerate}
  \item 异常在可见通道上严格为零；
  \item 块 Toeplitz--PSD 层级 \(\mathbf{T}^{\mathcal{A}_m}_N(F_m)\succeq 0\) 对一切 \(N\) 成立，
\end{enumerate}
则存在 \(\mathcal{A}_m\)-值正测度 \(\boldsymbol\mu_m\) 使 \(F_m\) 具有算子值 Herglotz 表示；
从而所有通道 \(F_{\overline\chi}\) 同时为同一 \(\boldsymbol\mu_m\) 的 Fourier--Pontryagin 投影（完全正可见实现存在）。
反之，若某群商 \(G\twoheadrightarrow Q\) 上存在上述完全正可见实现并使闭环接口严格交换，则该群商必唯一因子化通过 \(G^{\mathrm{BC}}_m\)；
因此 strictification 具有最大可见商的普适性质。
\end{conclusion}
\begin{proof}[证明]
见定理 \ref{thm:xi-gbc-maximal-cp-visible-realization}（并结合命题 \ref{prop:pom-bc-quotient-strictify}、\ref{prop:pom-bc-quotient-universal}）。证毕。
\end{proof}

\begin{conclusion}[端点阈值承载的中心性与通道强度谱：\texorpdfstring{$\boldsymbol\mu_m(\{-1\})$}{mu(-1)} 诱导不可约通道的非负标量分解]\label{con:xi-terminal-endpoint-atom-channel-strength-spectrum}
在结论 \ref{con:xi-terminal-strictification-maximal-cp-visible-quotient} 的完全正实现下，端点 \(-1\) 的原子
\(
\mathbf{M}_{-1}:=\boldsymbol\mu_m(\{-1\})\in(\mathcal{A}_m)_+
\)
存在且为正元。
若 \(G^{\mathrm{BC}}_m\) 为阿贝尔群，则 \(\mathcal{A}_m\) 交换从而 \(\mathbf{M}_{-1}\) 属于中心；
更一般地，若 \(\boldsymbol\mu_m\) 取值于中心子代数，则对任意不可约分量 \(\pi\) 有
\[
(\mathbf{M}_{-1})_\pi=\lambda_\pi\,I_{d_\pi},
\qquad
\lambda_\pi\ge 0.
\]
族 \(\{\lambda_\pi\}\) 给出端点阈值承载的有限维“通道强度分解”，把极端共振边界层强度压缩为不可约表示索引下的一组非负谱数据。
\end{conclusion}
\begin{proof}[证明]
见定理 \ref{thm:xi-endpoint-atom-central-channel-strength}。证毕。
\end{proof}

\begin{conclusion}[端点热核探针的严格恒等式：有限 Toeplitz 二次型单调抽取端点原子质量]\label{con:xi-terminal-endpoint-heat-kernel-probe-identity}
令 \(\mathscr{C}(w)=c_0+2\sum_{n\ge 1}c_n w^n\) 在 \(w=0\) 邻域解析，并约定 \(c_{-n}:=\overline{c_n}\)。
对每个 \(N\ge 0\) 定义 Toeplitz 主子矩阵 \(\mathbf{T}_N:=(c_{j-k})_{0\le j,k\le N}\)，并令二项式向量 \(b^{(N)}\in\RR^{N+1}\) 满足
\[
b^{(N)}_k:=(-1)^k\binom{N}{k}\qquad(0\le k\le N),
\qquad
a_N:=\frac{(b^{(N)})^{\mathsf T}\mathbf{T}_N b^{(N)}}{4^N}.
\]
若 \(\mathscr{C}\) 属于 Carath\'eodory 类，则存在唯一有限正测度 \(\mu\) 使 \(c_n=\int_{\TT}\xi^{-n}\,\dd\mu(\xi)\)，并且成立严格恒等式
\[
a_N=\int_{\TT}\left|\frac{1-\xi}{2}\right|^{2N}\,\dd\mu(\xi)\ge 0,
\qquad
a_N\downarrow \mu(\{-1\})\quad(N\to\infty).
\]
因此端点阈值承载 \(\mu(\{-1\})\) 可由有限 Toeplitz 数据以单调方式无歧义提取。
\end{conclusion}
\begin{proof}[证明]
见定理 \ref{thm:xi-endpoint-heat-kernel-probe}。证毕。
\end{proof}

\begin{conclusion}[端点原子提取的指数型误差界：支撑受控时达到指数逼近]\label{con:xi-terminal-endpoint-atom-exponential-error}
在结论 \ref{con:xi-terminal-endpoint-heat-kernel-probe-identity} 的设定下，取任意闭集 \(K\subset\TT\setminus\{-1\}\)，并令
\[
q(K):=\sup_{\xi\in K}\left|\frac{1-\xi}{2}\right|\in(0,1).
\]
则对所有 \(N\ge 0\) 有
\[
0\le a_N-\mu(\{-1\})\le q(K)^{2N}\mu(K)+\mu\!\bigl(\TT\setminus(\{-1\}\cup K)\bigr).
\]
特别地，若 \(\mu(\TT\setminus(\{-1\}\cup K))=0\)，则
\(
0\le a_N-\mu(\{-1\})\le q(K)^{2N}\mu(K)
\)，从而在端点邻域之外存在谱隙（或支撑受控）时，端点原子质量由 \(\{a_N\}\) 以指数速度逼近提取；其中 \(N\) 充当严格的端点分辨率深度。
\end{conclusion}
\begin{proof}[证明]
见推论 \ref{cor:xi-endpoint-heat-kernel-exp-error}。证毕。
\end{proof}

\begin{conclusion}[阈值强度与端点谱维数：余项衰减指数精确等价于端点局部幂律指数]\label{con:xi-terminal-endpoint-threshold-spectral-dimension}
将表示测度分解为
\[
\mu=\mathfrak m_{-1}\delta_{-1}+\mu_{\neq -1},
\qquad
\mathfrak m_{-1}:=\mu(\{-1\})\ge 0.
\]
则对每个 \(N\ge 0\) 有严格分离恒等式
\[
a_N=\mathfrak m_{-1}+\int_{\TT\setminus\{-1\}}\left|\frac{1-\xi}{2}\right|^{2N}\,\dd\mu_{\neq-1}(\xi),
\qquad
a_N-\mathfrak m_{-1}\downarrow 0,
\]
从而 \(\mathfrak m_{-1}\) 为端点阈值承载的唯一零阶强度参数。进一步，若 \(\mu_{\neq -1}\) 在端点角窗满足局部幂律：存在 \(\alpha>-1\) 与常数 \(K_\pm\ge 0\) 使当 \(\varepsilon\downarrow 0\) 时
\[
\mu_{\neq -1}\bigl(\{\pi<\arg\xi<\pi+\varepsilon\}\bigr)=K_+\,\varepsilon^{\alpha+1}+o(\varepsilon^{\alpha+1}),
\qquad
\mu_{\neq -1}\bigl(\{\pi-\varepsilon<\arg\xi<\pi\}\bigr)=K_-\,\varepsilon^{\alpha+1}+o(\varepsilon^{\alpha+1}),
\]
则余项具有幂律渐近
\[
a_N-\mathfrak m_{-1}\sim (\alpha+1)\,2^{\alpha}\,
\Gamma\!\left(\frac{\alpha+1}{2}\right)\,
(K_++K_-)\,N^{-(\alpha+1)/2}\qquad (N\to\infty).
\]
尤其当 \(\mu\) 在 \(\xi=-1\) 邻域绝对连续且角密度 \(f(\theta)\) 于 \(\theta=\pi\) 连续且 \(f(\pi)>0\) 时，\(\alpha=0\) 且
\[
a_N-\mathfrak m_{-1}\sim 2f(\pi)\sqrt{\frac{\pi}{N}}.
\]
故端点局部谱维数指数（定义 \ref{def:xi-endpoint-spectral-dimension}）可由 \(a_N-\mathfrak m_{-1}\) 的幂律衰减率直接识别。
\end{conclusion}
\begin{proof}[证明]
见推论 \ref{cor:xi-endpoint-atom-separation} 与定理 \ref{thm:xi-endpoint-spectral-dimension-laplace}。证毕。
\end{proof}

\begin{conclusion}[Adams--Toeplitz--PSD 派生层级：Carath\'eodory 正性在幂覆盖下闭合]\label{con:xi-terminal-adams-toeplitz-psd-hierarchy}
若 \(\mathscr{C}(w)=c_0+2\sum_{n\ge 1}c_n w^n\) 属于 Carath\'eodory 类并由正测度 \(\mu\) 表示，则对每个 \(m\ge 1\) 定义推前测度 \(\mu_m:=(\xi\mapsto\xi^m)_\ast\mu\)。其矩满足
\[
c^{(m)}_n:=\int_{\TT}\xi^{-n}\,\dd\mu_m(\xi)=c_{mn}.
\]
因此对任意 \(N,m\ge 1\)，稀疏 Toeplitz 主子矩阵族
\[
\mathbf{T}^{(m)}_N:=\bigl(c_{m(j-k)}\bigr)_{0\le j,k\le N}\succeq 0
\]
自动成立。该派生层级在不改变正性证书口径的前提下，将端点压缩与外置尺度参数 \(m\) 直接耦合，形成一族可随 \(m\) 调节分辨率的多尺度 PSD 约束。
\end{conclusion}
\begin{proof}[证明]
见命题 \ref{prop:xi-adams-pushforward-toeplitz-principal-minor-monotonicity} 与式 \eqref{eq:xi-adams-toeplitz-psd}。证毕。
\end{proof}

\begin{conclusion}[尺度交换律：幂覆盖将端点边界层缺陷搬运到任意紧子圆盘并产生高度无关的 Jensen 正缺陷]\label{con:xi-terminal-scale-exchange-power-covering-jensen}
在自对偶完成化与固定切片制度下，离切片点状缺陷的圆盘像 \(w_\rho\) 由 Cayley--Busemann 结构性压缩到端点边界层（定理 \ref{thm:xi-offcritical-quadratic-radial-compression}）。
沿幂覆盖 \(\pi_m(\omega)=\omega^m\) 的 Adams--Norm 折叠门可将该端点边界层中的点状缺陷搬运到任意预设的紧子圆盘：对任意 \(r_\star\in(0,1)\) 存在显式尺度 \(m_\star\) 使 \(|w_\rho|^{m_\star}\le r_\star\)。
并且在固定半径 \(R\in(r_\star,1)\) 处产生高度无关的 Jensen 正缺陷下界：
\[
D_{\mathrm{Nm}_{m_\star}[F_\zeta]}(R)\ \ge\ \log\!\frac{R}{r_\star}\ >\ 0.
\]
该结论给出一条制度内可审计的交换机制：二次压缩把离切片信息压入端点边界层，而幂覆盖/折叠门以尺度参数 \(m_\star\) 将其解压到固定半径并以 Jensen 缺陷显影。
\end{conclusion}
\begin{proof}[证明]
见推论 \ref{cor:xi-adams-norm-fixed-radius-jensen-certificate}（并结合 \eqref{eq:xi-adams-norm-mstar-def}--\eqref{eq:xi-adams-norm-mstar-radius}）。证毕。
\end{proof}

\begin{conclusion}[端点重整化的二元正交分解：普适 \texorpdfstring{$\zeta$}{zeta}-极 \texorpdfstring{$\oplus$}{oplus} 离散 \texorpdfstring{$\chi_4$}{chi4} 本征残余]\label{con:xi-terminal-endpoint-renormalization-zeta-chi4-orthogonal}
在完成化坐标 \(S=r+r^{-1}\) 下，最小反变纤维坐标为 \(\delta:=r-r^{-1}\)，并满足反演反变部分的最小闭包
\(
\mathcal{A}^{-}=\delta\,\ZZ[S]
\)
（备注 \ref{rem:endpoint-inv-antiinv-min-closure}）。
奇 Chebyshev--Adams 伴随族满足恒等式
\[
D_d(S,\delta)=\delta\,U_{d-1}(S/2),
\]
从而在端点分支 \(r=\pm\mathrm{i}\) 上 \(S=0\) 并有
\[
\frac{D_d(0,\delta)}{\delta}=U_{d-1}(0)=\sin(d\pi/2)=:\varepsilon(d)\in\{-1,0,1\}.
\]
该端点通道指标 \(\varepsilon(d)\) 给出模 \(4\) 的完整奇类分裂，并与非平凡 Dirichlet 角色 \(\chi_4\) 严格一致：\(\varepsilon(d)=\chi_4(d)\)（定理 \ref{thm:endpoint-four-channel-chi4}）。
进一步，反变端点观测在 \(\sigma\downarrow 0\) 处解析并满足
\[
\lim_{\sigma\downarrow 0}\mathcal O(\mathcal W_{\sigma}^{\delta}G)
=L(1,\chi_4)\,\mathcal O(G)=\frac{\pi}{4}\,\mathcal O(G),
\]
而与之正交地，偶通道的全部发散仅沿 rank-one 方向出现并由 \(\zeta(1+\sigma)\) 的极点完全控制（定理 \ref{thm:endpoint-even-finite-part-renormalization}）。
因此端点重整化后的非紧性可被压缩为“普适 \(\zeta\)-极 \(\oplus\) 离散 \(\chi_4\) 本征残余”的二元正交分解。
\end{conclusion}
\begin{proof}[证明]
Chebyshev--Adams 恒等式与 \(\varepsilon=\chi_4\) 见命题 \ref{prop:endpoint-odd-chebyshev-identity} 与定理 \ref{thm:endpoint-four-channel-chi4}；
反变通道极限见推论 \ref{cor:endpoint-odd-analytic-at-zero}；
偶通道的极点控制与残余分解见注 \ref{rem:endpoint-renormalization-universal-pole-discrete-hair}（并参见注 \ref{rem:endpoint-obstruction-dichotomy}）。
证毕。
\end{proof}

\begin{conclusion}[有限阿贝尔可见商的块 Toeplitz 对角化：单一块 PSD 证书等价于全通道 PSD]\label{con:xi-terminal-abelian-block-toeplitz-diagonalization}
令 \(\mathcal{A}\) 为有限维 \(C^\ast\)-代数，并取解析函数 \(F:\mathbb{D}\to\mathcal{A}\)。
则以下命题等价：\(\Re F(w)\succeq 0\)（算子值 Carath\'eodory）、对每个不可约表示 \(\pi\) 有 \(\Re F_\pi(w)\succeq 0\)、以及对每个 \(N\ge 0\) 的块 Toeplitz 截断 \(\mathbf{T}^{\mathcal{A}}_N(F)\succeq 0\)。
特别地，当 \(\mathcal{A}=\CC[H]\) 且 \(H\) 为有限阿贝尔群时，在 Fourier--Pontryagin 对偶分解下 \(\mathbf{T}^{\mathcal{A}}_N(F)\) 与直和矩阵 \(\bigoplus_{\chi\in\widehat H}\mathbf{T}_N(F_\chi)\) 酉等价；因此
\[
\mathbf{T}^{\mathcal{A}}_N(F)\succeq 0\ \Longleftrightarrow\ \forall\chi\in\widehat H,\ \mathbf{T}_N(F_\chi)\succeq 0.
\]
从而“逐通道 PSD 核验”可被统一为一个单一块 Toeplitz--PSD 可行域（单一半正定证书）；与结论 \ref{con:xi-terminal-adams-toeplitz-psd-hierarchy} 的 Adams--Toeplitz 派生层级结合，可在单一块证书口径下实现全尺度、全通道的正性核验。
\end{conclusion}
\begin{proof}[证明]
见定理 \ref{thm:xi-fourier-pontryagin-block-toeplitz} 与推论 \ref{cor:xi-abelian-block-toeplitz-directsum}。证毕。
\end{proof}

\begin{conclusion}[异常类与散射绕数的同一阻碍：交换失败残差与拓扑绕数在可见口径下严格对齐]\label{con:xi-terminal-anom-winding-obstruction-equivalence}
固定分辨率 \(m\) 并考虑可见扭曲通道诱导的 unitary-slice 散射行列式回路
\(\{\det\mathcal S_{\overline\chi}\}_{\overline\chi\in\widehat{G^{\mathrm{BC}}_m}}\)。
若异常在可见通道上严格为零，则每个可见通道的散射回路可取高能归一化
\(\det\mathcal S_{\overline\chi}(\infty)=1\)，并满足
\[
\mathrm{Wind}\bigl(\det\mathcal S_{\overline\chi}\bigr)=0
\qquad(\forall\,\overline\chi\in\widehat{G^{\mathrm{BC}}_m}),
\]
从而在有限型点状共振口径下所有可见通道的真共振指数归零。
反向地，若对某群商 \(G\twoheadrightarrow H\) 的全部可见通道均可实现上述归一化与零绕数，则该群商必唯一因子化通过 \(G^{\mathrm{BC}}_m\)，并在 \(H\) 上使异常严格为零。
因此“代数残差”（异常类）与“拓扑绕数”在可见口径下构成同一阻碍类。
\end{conclusion}
\begin{proof}[证明]
见定理 \ref{thm:xi-anom-winding-obstruction}（真共振指数与绕数/Blaschke 度的对齐由定理 \ref{thm:selfdual-wind-deg-dbl-kappa-identity} 给出）。证毕。
\end{proof}

\begin{conclusion}[有界切片的有限前沿原理：证书搜索与编译优化在每一分辨率层级均可退化为有限枚举]\label{con:xi-terminal-finite-frontier-principle}
在任意固定有界复杂度切片 \(\mathbf{PW}^{\ast}_{\le}\) 内，取保守语义函子 \(\mathsf{Sem}\) 与任一代价函子 \(\mathcal{C}\)。
则对任意投影词 \(w\in \mathbf{PW}^{\ast}_{\le}\)，其语义等价类 \([w]\) 在代价域中的像 \(\mathcal{C}([w])\subseteq \mathbb{T}^{(2)}\) 具有有限帕累托前沿，并可由有限枚举构造。
因此任何以“接口正规形参数—异常签名—谱指纹”为核对输出的证书目标，在每一固定分辨率层级上必然退化为有限枚举与有限核对；困难被严格集中到“切片上界随分辨率增长律”而非实现细节。
\end{conclusion}
\begin{proof}[证明]
见定理 \ref{thm:xi-finite-frontier-certificate}。证毕。
\end{proof}

\begin{conclusion}[输入合法性约束的几何化产生常数级谱跃迁：加法投影碰撞常数在真实核上显著下降]\label{con:xi-terminal-addition-collision-spectrum-constant-factor-drop}
把“加法作为一次投影”的同像碰撞谱常数记为 \(r^{\mathrm{add}}_q\)（为相应 \(q\)-碰撞 moment-kernel 的 Perron 根；见附录 \ref{app:add-collision-spectrum}）。
无约束输入（10 状态归一化核驱动）与真实输入约束（40 状态核编码 Zeckendorf 合法性）之间存在显著常数差异：
\[
r^{\mathrm{add}}_{2,\mathrm{sync}}\approx 8.7430515,\qquad
r^{\mathrm{add}}_{2,\mathrm{real}}\approx 3.9985433,
\qquad
\frac{r^{\mathrm{add}}_{2,\mathrm{sync}}}{r^{\mathrm{add}}_{2,\mathrm{real}}}\approx 2.186559.
\]
该“指数尺度的常数因子级下降”提供一个可直接报告的价值锚点：真实输入核不仅更贴近算术对象，而且在空间损失谱（以碰撞主尺度表征）上携带不可由无约束归一化器替代的独立信息。
\end{conclusion}
\begin{proof}[证明]
见命题 \ref{prop:addition-collision-spectrum-sync10-vs-real40} 与表 \ref{tab:add-collision-spectrum-10-vs-40}。证毕。
\end{proof}

\begin{conclusion}[折叠输出的内生 Gibbs 规格与 Parry 基线几乎同熵：最大熵缺口降至微比特级且 KL 率等同于熵缺口]\label{con:xi-terminal-fold-output-gibbs-parry-microbit-gap}
对 Fold 输出分布 \(\pi_m(x)=d_m(x)/2^m\) 的块频率拟合得到两状态 Markov 规格 \(P_{\mathrm{fit}}\)，其一步转移与黄金均值 Parry 基线仅在 \(10^{-3}\) 量级偏离（例如 \(P_{\mathrm{fit}}(1\mid 0)=0.3829949\) 对比 \(P_{\mathrm{Parry}}(1\mid 0)=0.3819660\)）。
更关键的是熵率层面的微缺口：
\[
h_{\mathrm{fit}}\approx 0.481210204266,\quad
\log\varphi\approx 0.481211825060,\quad
\Delta h:=\log\varphi-h_{\mathrm{fit}}\approx 1.620794\times 10^{-6}\ \text{nats/symbol}.
\]
并且 KL 散度率满足严格一致性恒等式
\[
D_{\mathrm{KL}}^{\mathrm{rate}}(P_{\mathrm{fit}}\|P_{\mathrm{Parry}})=\log\varphi-h_{\mathrm{fit}}=\Delta h,
\]
从而折叠诱导的“偏离”主要体现为常数级指纹（有限部/端点层），而非熵率层面的结构崩塌。
\end{conclusion}
\begin{proof}[证明]
见注 \ref{rem:pom-gibbs-markov-fit} 与注 \ref{rem:pom-gibbs-parry-entropy-gap}（其中数值证书由 \texttt{scripts/exp\_fold\_output\_gibbs\_markov\_fit.py} 生成并以 \texttt{\string\input} 方式纳入正文）。证毕。
\end{proof}

\begin{conclusion}[端点质量给出黄金常数的过定约束审计：两条独立比值估计在 Parry 模型下必同收敛于 \texorpdfstring{$\varphi$}{phi}]\label{con:xi-terminal-parry-endpoint-mass-overdetermined-phi-audit}
以任意经验直方图 \(\hat\pi_{m,N}\) 定义端点质量
\[
\widehat M_{ij}(m):=\sum_{\substack{u\in X_m\\ u_1=i,\ u_m=j}}\hat\pi_{m,N}(u).
\]
若数据在分辨率 \(m\) 下接近黄金均值 Parry 基线，则可构造两条相互独立的黄金常数估计量：
\[
\widehat\varphi_{(00/01)}(m):=\frac{\widehat M_{00}(m)}{\widehat M_{01}(m)}\cdot \frac{F_{m-1}}{F_m},
\qquad
\widehat\varphi_{(01/11)}(m):=\frac{\widehat M_{01}(m)}{\widehat M_{11}(m)}\cdot \frac{F_{m-2}}{F_{m-1}}.
\]
在 Parry 真模型下，两者在极限意义下收敛到同一常数 \(\varphi\)，从而形成一个低维、过定约束的自洽审计机制：任何偏离 Parry 的端点结构将使两条比值估计产生可检测的不一致。
\end{conclusion}
\begin{proof}[证明]
见推论 \ref{cor:parry-endpoint-phi-estimators}。证毕。
\end{proof}

\begin{conclusion}[模群判据下的贝叶斯逆与零知识因子化：完全正信道的“可见闭包不变性”给出充要刻画]\label{con:xi-terminal-bayesian-inverse-zk-factorization-invariance}
在算子代数设定中，贝叶斯逆 \(F^\sharp\) 的存在性由 Tomita--Takesaki 模群判据控制。
一旦 \(F^\sharp\) 存在，则对任意正规完全正转录信道 \(\Phi:\mathcal{M}\to\mathcal{T}\)，以下命题等价：
\begin{enumerate}
  \item 存在正规完全正映射 \(S:\mathcal{N}\to\mathcal{T}\) 使 \(\Phi=S\circ F^\sharp\)（相对于 \(\mathcal{N}\) 的完美零知识因子化）；
  \item \(\Phi\) 在可见闭包幂等子 \(K:=F\circ F^\sharp\) 下不变：\(\Phi=\Phi\circ K\)。
\end{enumerate}
并且在等价条件下分解 \(\Phi=S\circ F^\sharp\) 中的 \(S\) 唯一且满足 \(S=\Phi|_{\mathcal{N}}\)。
因此“零知识/可因子化”从构造性问题降格为单一不变性检验。
\end{conclusion}
\begin{proof}[证明]
见定理 \ref{thm:op-algebra-bayes-inverse-zk}。证毕。
\end{proof}

